% ---------------------------------------------------------------------------
% Preámbulo común del libro. Se inyecta con  pandoc -H
% Define la identidad visual: cajas pedagógicas, colores y microtipografía.
% ---------------------------------------------------------------------------
\usepackage{tcolorbox}
\tcbuselibrary{skins,breakable}
\usepackage{amsmath,amssymb}
\usepackage{booktabs}
\usepackage{microtype}

% --- Símbolos de dificultad y de nivel -------------------------------------
% Latin Modern no tiene glifos para ○ ◐ ● ★ ◆, y XeTeX los descarta EN SILENCIO:
% el marcador de dificultad de cada problema desaparecía del PDF sin aviso.
% Se mapean a equivalentes matemáticos que sí existen.
\usepackage[nointegrals]{wasysym}
\usepackage{newunicodechar}
\newunicodechar{○}{\ensuremath{\Circle}}
\newunicodechar{◐}{\ensuremath{\LEFTcircle}}
\newunicodechar{●}{\ensuremath{\CIRCLE}}
% \bigstar y \blacklozenge NO valen aquí: la plantilla de pandoc carga
% unicode-math, que los reasigna a caracteres que la fuente matemática no tiene.
% Los dingbats de pifont son independientes de la fuente matemática.
\usepackage{pifont}
\newunicodechar{★}{\ding{72}}
\newunicodechar{◆}{\ding{117}}

% --- Fuente monoespaciada con cobertura Unicode ----------------------------
% Los diagramas ASCII y las plantillas del libro usan caracteres de dibujo de
% cajas (─ │ ├ └ ┬) y casillas (□). Latin Modern Mono no los tiene y XeTeX los
% descarta EN SILENCIO: los diagramas salían desmontados sin ningún aviso.
% DejaVu Sans Mono sí los cubre. Si no está instalada, se sigue sin ella.
\IfFontExistsTF{DejaVu Sans Mono}{%
  \setmonofont{DejaVu Sans Mono}[Scale=MatchLowercase]%
}{}

% Los mismos caracteres, cuando aparecen fuera de un bloque de código.
% Los cuadrados de Zapf Dingbats llevan sombra; para una lista que se imprime
% y se marca conviene una caja limpia, dibujada y por tanto independiente de la
% fuente.
\newcommand{\casilla}{\raisebox{-0.08em}{{\setlength{\fboxsep}{0pt}%
  \framebox[0.72em]{\rule{0pt}{0.72em}}}}}
\newunicodechar{□}{\casilla}
\newunicodechar{✓}{\ding{51}}
\newunicodechar{⁻}{\ensuremath{^{-}}}
\newunicodechar{⁰}{\ensuremath{^{0}}}
\newunicodechar{⁴}{\ensuremath{^{4}}}
\newunicodechar{⁵}{\ensuremath{^{5}}}
\newunicodechar{⁶}{\ensuremath{^{6}}}
\newunicodechar{₀}{\ensuremath{_{0}}}
\newunicodechar{₂}{\ensuremath{_{2}}}
\newunicodechar{ᵢ}{\ensuremath{_{i}}}
\newunicodechar{ⱼ}{\ensuremath{_{j}}}
\newunicodechar{ṁ}{\ensuremath{\dot m}}

% Letras griegas y símbolos que aparecen en el texto corrido, fuera de modo
% matemático. Latin Modern Roman no las tiene y también se caían en silencio.
% Nada de upgreek: choca con unicode-math + microtype a través de mt-eur.cfg.
\newunicodechar{μ}{\ensuremath{\mu}}
\newunicodechar{σ}{\ensuremath{\sigma}}
\newunicodechar{π}{\ensuremath{\pi}}
\newunicodechar{ρ}{\ensuremath{\rho}}
\newunicodechar{α}{\ensuremath{\alpha}}
\newunicodechar{τ}{\ensuremath{\tau}}
\newunicodechar{χ}{\ensuremath{\chi}}
\newunicodechar{ε}{\ensuremath{\varepsilon}}
\newunicodechar{γ}{\ensuremath{\gamma}}
\newunicodechar{Δ}{\ensuremath{\Delta}}
\newunicodechar{Σ}{\ensuremath{\Sigma}}
\newunicodechar{∝}{\ensuremath{\propto}}
\newunicodechar{∞}{\ensuremath{\infty}}
\newunicodechar{≪}{\ensuremath{\ll}}
\newunicodechar{∼}{\ensuremath{\sim}}
\newunicodechar{⟨}{\ensuremath{\langle}}
\newunicodechar{⟩}{\ensuremath{\rangle}}

% Operadores que también aparecen sueltos en el texto corrido.
\newunicodechar{≈}{\ensuremath{\approx}}
\newunicodechar{≥}{\ensuremath{\geq}}
\newunicodechar{≤}{\ensuremath{\leq}}
\newunicodechar{÷}{\ensuremath{\div}}
\newunicodechar{−}{\ensuremath{-}}
\newunicodechar{√}{\ensuremath{\surd}}
\newunicodechar{≠}{\ensuremath{\neq}}

% Los caracteres de dibujo de cajas sólo aparecen dentro de bloques literales,
% donde los cubre DejaVu Sans Mono. Por si alguna vez aparecen en texto corrido,
% se fuerza la fuente monoespaciada. \symbol no vuelve a pasar por newunicodechar,
% así que no hay recursión.
\newunicodechar{─}{{\ttfamily\symbol{"2500}}}
\newunicodechar{│}{{\ttfamily\symbol{"2502}}}
\newunicodechar{┌}{{\ttfamily\symbol{"250C}}}
\newunicodechar{┐}{{\ttfamily\symbol{"2510}}}
\newunicodechar{└}{{\ttfamily\symbol{"2514}}}
\newunicodechar{┘}{{\ttfamily\symbol{"2518}}}
\newunicodechar{├}{{\ttfamily\symbol{"251C}}}
\newunicodechar{┤}{{\ttfamily\symbol{"2524}}}
\newunicodechar{┬}{{\ttfamily\symbol{"252C}}}
\newunicodechar{┴}{{\ttfamily\symbol{"2534}}}
\newunicodechar{┼}{{\ttfamily\symbol{"253C}}}

% El interlineado del cuerpo separa las líneas de los diagramas ASCII y rompe
% los trazos verticales. Dentro de un bloque literal se vuelve al natural.
\usepackage{etoolbox}
\AtBeginEnvironment{verbatim}{\setstretch{1}}
\AtBeginEnvironment{Shaded}{\setstretch{1}}

% --- Paleta (idéntica a herramientas/estilo_libro.py) ----------------------
\definecolor{tinta}{HTML}{1B2A41}
\definecolor{azul}{HTML}{2F6FA8}
\definecolor{rojo}{HTML}{C1443C}
\definecolor{verde}{HTML}{3F8F6B}
\definecolor{ocre}{HTML}{D99A2B}
\definecolor{morado}{HTML}{7D5BA6}
\definecolor{tealc}{HTML}{3A8FA0}
\definecolor{grisc}{HTML}{8A94A6}
\definecolor{papel}{HTML}{FBFAF7}
\definecolor{claro}{HTML}{E8ECF2}

% hyperref lo carga la plantilla de pandoc DESPUÉS de este bloque:
% por eso los colores de enlace se fijan al empezar el documento.
\AtBeginDocument{\hypersetup{colorlinks=true,linkcolor=azul,urlcolor=azul,citecolor=verde}}

% --- Estilo base de todas las cajas ---------------------------------------
\tcbset{
  cajalibro/.style={
    enhanced, breakable, sharp corners=downhill,
    boxrule=0pt, leftrule=2.2pt, toprule=0pt, bottomrule=0pt, rightrule=0pt,
    left=8pt, right=8pt, top=7pt, bottom=7pt,
  }
}

% Rótulo de cabecera de caja. Se inserta como contenido, no como "title":
% el mecanismo de títulos de tcolorbox no sobrevive bien a los cortes de
% página de las cajas breakable.
\newcommand{\cajatitulo}[2]{%
  {\bfseries\sffamily\small\textcolor{#1}{#2}}\par\smallskip}

\newtcolorbox{cajapregunta}{cajalibro, colback=azul!6, colframe=azul,
  before upper={\cajatitulo{azul}{UNA PREGUNTA}}}
\newtcolorbox{cajaantes}{cajalibro, colback=ocre!8, colframe=ocre,
  before upper={\cajatitulo{ocre!80!black}{ANTES DE CALCULAR}}}
\newtcolorbox{cajaherramientas}{cajalibro, colback=claro, colframe=tinta,
  before upper={\cajatitulo{tinta}{CAJA DE HERRAMIENTAS MATEMÁTICA}}}
\newtcolorbox{cajahistoria}{cajalibro, colback=verde!6, colframe=verde,
  before upper={\cajatitulo{verde!60!black}{HISTORIA}}}
\newtcolorbox{cajaexplica}{cajalibro, colback=morado!6, colframe=morado,
  before upper={\cajatitulo{morado}{EXPLÍCALO SIN ESCONDERTE DETRÁS DE LAS ECUACIONES}}}
\newtcolorbox{cajaesencial}{cajalibro, colback=tinta!6, colframe=tinta,
  before upper={\cajatitulo{tinta}{LO ESENCIAL}}}
\newtcolorbox{cajaabierto}{cajalibro, colback=tealc!7, colframe=tealc,
  before upper={\cajatitulo{tealc!70!black}{PREGUNTAS QUE QUEDAN ABIERTAS}}}
\newtcolorbox{cajaexperimento}{cajalibro, colback=azul!5, colframe=azul!70!black,
  before upper={\cajatitulo{azul!70!black}{EXPERIMENTO COMPUTACIONAL}}}
\newtcolorbox{cajaaviso}{cajalibro, colback=rojo!6, colframe=rojo,
  before upper={\cajatitulo{rojo}{TRAMPA}}}
\newtcolorbox{cajasupuestos}{cajalibro, colback=grisc!10, colframe=grisc,
  before upper={\cajatitulo{tinta}{¿QUÉ ESTAMOS SUPONIENDO?}}}
\newtcolorbox{cajafalla}{cajalibro, colback=rojo!5, colframe=rojo!70!black,
  before upper={\cajatitulo{rojo!70!black}{¿CUÁNDO FALLA?}}}
\newtcolorbox{cajaia}{cajalibro, colback=morado!5, colframe=morado!70!black,
  before upper={\cajatitulo{morado!70!black}{PROTOCOLO CON IA}}}
\newtcolorbox{cajajuega}{cajalibro, colback=ocre!6, colframe=ocre!80!black,
  before upper={\cajatitulo{ocre!60!black}{JUEGA CON EL MODELO}}}
\newtcolorbox{cajanumero}{cajalibro, colback=claro, colframe=grisc,
  before upper={\cajatitulo{tinta}{NÚMEROS QUE CONVIENE SABERSE}}}

% Cita destacada, sin caja
\newenvironment{citadestacada}
  {\begin{quote}\itshape\color{tinta!85}}
  {\end{quote}}

% Figuras: pie en gris y pequeño
\usepackage{caption}
\captionsetup{font=small, labelfont={bf,color=tinta}, textfont={color=tinta!80}}
